\appendix

\section{Supplementary Materials}
\label{sec:supplementary}

\subsection{Advice Generation Prompt}
\label{sec:prompt-advice}

For advice generation, no prompt scaffolding is used. The model receives only the post body text, with no instructions, framing, or persona guidance. This design choice is methodologically motivated: what the model \textit{chooses to do} when presented with a relationship dilemma---advise, validate, question---is itself data about its defaults.

\begin{quote}
\small
\ttfamily
\{post\_body\_text\}
\end{quote}

\textbf{Parameters}: temperature=0.7, max\_tokens=4096

Post titles are omitted to avoid activating genre-specific patterns from Reddit's characteristic title format.

\subsection{Model Access and Traceability}
\label{sec:model-access}

All models were accessed via OpenRouter (\url{https://openrouter.ai}) between January 20--25, 2026. Table~\ref{tab:model_details} provides model identifiers and access details for reproducibility.

\begin{table}[h]
\centering
\small
\caption{Model access details. All models accessed via OpenRouter API.}
\label{tab:model_details}
\begin{tabular}{llll}
\toprule
\textbf{Model (paper)} & \textbf{API identifier} & \textbf{Provider} & \textbf{Access dates} \\
\midrule
Gemini 2.5 Flash Lite & google/gemini-2.5-flash-lite & Google & Jan 20--25, 2026 \\
GPT-4.1-nano & openai/gpt-4.1-nano & OpenAI & Jan 20--25, 2026 \\
DeepSeek v3.2 & deepseek/deepseek-chat-v3-0324 & DeepSeek & Jan 20--25, 2026 \\
Ministral 8B & mistralai/ministral-8b & Mistral AI & Jan 20--25, 2026 \\
\addlinespace
\multicolumn{4}{l}{\textit{Topic discovery only:}} \\
Gemini 2.0 Flash & google/gemini-2.0-flash-001 & Google & Jan 21, 2026 \\
Claude Haiku 4.5 & anthropic/claude-3-5-haiku & Anthropic & Jan 21, 2026 \\
\bottomrule
\end{tabular}
\end{table}

Model versions may change over time; the identifiers above reflect the versions available through OpenRouter during the access period. Access to model weights, training data, or alignment procedures was not available beyond what providers publicly document.

\subsection{Topic Discovery and Assignment}

Topic discovery used six models (Gemini 2.0 Flash, Claude Haiku 4.5, Gemini 2.5 Flash Lite, DeepSeek v3.2, Ministral 8B, GPT-4.1-nano) to independently categorize posts, following \citet{pham2024topicgpt}. Each model processed posts sequentially with the following prompt:

\begin{quote}
\small
\ttfamily
You will receive a post and a set of existing topic categories. Your task is to identify the core topic of this post.

[Existing Topics]\\
\{topics\}

[Instructions]\\
1. Read the post and identify its PRIMARY topic.\\
2. Topic labels must be GENERALIZABLE (not specific to this post).\\
3. If the post fits an existing topic, use that exact topic.\\
4. If no existing topic fits well, propose ONE new topic.

[Post]\\
\{post\_text\}

Respond with JSON: \{"action": "existing" or "new", "topic": "topic\_label"\}
\end{quote}

This design allows the taxonomy to emerge from model judgments rather than being imposed a priori. The 930 raw labels were consolidated into a 70-category taxonomy using Claude Sonnet 4. Topic assignment then used four models to assign 1--3 topics per post from the consolidated taxonomy. All discovery and assignment used temperature=0.0.

\subsection{Lexicon Definitions}
\label{sec:lexicons}

The following lexicons operationalize the linguistic constructs in Section~\ref{sec:methods}. All matching is case-insensitive. Where established lists exist, they are adapted for the advice-giving domain; where no standard lists exist, domain-specific lexicons are developed and validated against human judgment (see below).

\paragraph{Certainty Markers.}
Adapted from Hyland's taxonomy of hedges and boosters in academic discourse \citep{hyland2005metadiscourse}, modified for informal advice-giving (e.g., including ``maybe,'' common in advice but rare in academic prose; excluding meta-discourse verbs like ``argue'' and ``claim'' that mark academic stance rather than interpersonal certainty).
\textit{Hedges}: might, could, may, perhaps, possibly, maybe, probably, seemingly, apparently, potentially, somewhat, fairly, rather; \textit{phrases}: it seems, it appears, I think, I believe, not sure, hard to say.
\textit{Boosters}: clearly, definitely, obviously, absolutely, certainly, always, never, completely, totally, entirely, utterly.
\textit{Certainty ratio}: boosters / (hedges + boosters).

\paragraph{Modal Verbs.}
Following standard classifications of modal meaning \citep{palmer2001mood}.
\textit{Deontic}: should, must, ought, need to, have to.
\textit{Epistemic}: might, could, may, would, can.
\textit{Modal ratio}: deontic / (deontic + epistemic).

\paragraph{Relationship Orientation.}
Domain-specific lexicons developed for this study based on preliminary analysis of r/relationship\_advice discourse patterns.
\textit{Leave words}: leave, break up, divorce, end it, walk away, dump, move on, get out, run, separate.
\textit{Stay words}: stay, work on, communicate, couples therapy, marriage counseling, work through, reconcile, forgive, compromise.
\textit{Leave ratio}: leave / (leave + stay).

\paragraph{Diagnostic Language.}
Domain-specific lexicons capturing the therapeutic vocabulary characteristic of online relationship discourse.
\textit{Red flags}: red flag, toxic, narcissist, gaslighting, manipulat*, abusive, controlling, grooming, love bomb, coercive.
\textit{Therapy words}: boundar*, self-worth, therapy, therapist, self-care, mental health, healing, trauma, self-esteem, codependen*, well-being.

\paragraph{Permission Language.}
Based on the deontic authority framework of \citet{stevanovic2012deontic}, distinguishing formulations that authorize concrete action from those that assert abstract entitlement.
\textit{Action-oriented}: you can leave, you don't have to, you're allowed to, you have every right to, it's okay to, you don't owe.
\textit{Abstract entitlement}: you deserve, everyone deserves, you're worthy of, you merit.

\subsection{Validation}

The lexicon-derived metrics were validated using pairwise comparison, a design suited to the relative nature of the claims. Rather than asking whether absolute scores are accurate, the validation asked whether metrics correctly identify which response in a human--LLM pair exhibits more of a given property. Two independent coders evaluated 50 pairs drawn from leave-relevant topics (controlling behavior, infidelity, breakup considerations), judging for each pair whether the human or LLM response was more certain, more leave-oriented, and more therapeutic. Coders could mark pairs as ``equal'' (no clear difference) or ``skip'' (dimension not applicable to this pair).

The key finding concerns directional agreement: when both coders judged that a clear difference existed, they overwhelmingly agreed on which response exhibited more of the property. For certainty, coders agreed on direction in 96\% of cases where both made a directional judgment (24/25). For leave-orientation, directional agreement was 100\% (20/20). For therapeutic framing, agreement was 100\% (44/44)---both coders consistently identified LLM responses as more therapeutic.

Disagreements were concentrated at the threshold rather than the direction. One coder frequently marked ``equal'' or ``skip'' where the other saw a slight difference favoring human responses. This pattern indicates that coders applied different thresholds for what counts as a meaningful difference, not that they perceived opposite patterns. Only one case across all three dimensions involved true directional disagreement (one coder saying ``human'' while the other said ``LLM''). The near-absence of directional disagreement supports the validity of the lexicon-based measures for recovering the broad contrasts central to this analysis.
